% Generated draft: requires mathematical and source review.
\documentclass[11pt]{article}
\usepackage[margin=25mm]{geometry}
\usepackage{fontspec}
\setmainfont{Latin Modern Roman}
\usepackage{amsmath,amssymb,mathtools,unicode-math}
\setmathfont{Latin Modern Math}
\usepackage{enumitem,needspace,xurl,hyperref}
\setlength{\parindent}{0pt}
\setlength{\parskip}{5pt}
\setlength{\emergencystretch}{3em}
\newcommand{\N}{\mathbb{N}}
\newcommand{\Z}{\mathbb{Z}}
\newcommand{\Q}{\mathbb{Q}}
\newcommand{\R}{\mathbb{R}}
\newcommand{\C}{\mathbb{C}}
\newcommand{\F}{\mathbb{F}}
\newcommand{\A}{\mathbb{A}}
\newcommand{\PP}{\mathbb{P}}
\newcommand{\E}{\mathbb{E}}
\newcommand{\eps}{\varepsilon}
\newcommand{\dd}{\,\mathrm{d}}
\begin{document}
\section*{UnsolvedMath ID: 30001406}
\textbf{Source problem code:} catalogue entry
\subsection{Definitions and mathematical statement}
For terminology, consider the usual setting of a complete first-order theory \(T\). A formula \(\varphi(x;y)\), with finite tuples of variables \(x,y\), has the independence property if, for every positive integer \(n\), there are a model \(M\models T\), tuples \(a_1,\ldots,a_n\), and tuples \(b_I\) for all \(I\subseteq\{1,\ldots,n\}\), of the appropriate sorts, such that
\[
M\models\varphi(a_i;b_I)\quad\Longleftrightarrow\quad i\in I.
\]
The theory \(T\) is NIP if no formula has the independence property. For a fixed finite tuple of variables \(x\), \(S_x(M)\) denotes the set of complete types over \(M\) in those variables. These standard conventions provide terminology; the supplied record does not explicitly specify the class of theories or the meaning of its unindexed notation \(S(M)\).

The source record describes the following model-cardinality property for NIP theories:
\[
\exists M\models T\qquad |M|=\kappa\quad\text{and}\quad |S(M)|=\kappa^{|T|},
\]
possibly with the additional requirement that \(M\) be ``gross.'' Here the intended range and quantification of the cardinal \(\kappa\), the convention for \(|T|\), the type space denoted by \(S(M)\), and the definition of ``gross'' are not supplied. Consequently this display records the supplied expression, not a fully specified assertion.

The requested problem is to find a strengthening of this property that characterizes NIP. Schematically, one seeks a condition \(\mathcal P(T)\) such that
\[
\mathcal P(T)\quad\Longleftrightarrow\quad T\text{ is NIP},
\]
where \(\mathcal P\) strengthens the intended model-cardinality property. Neither the strengthening nor the permissible form of such a condition is specified. Thus this is an incompletely specified research request, not a recovered precise conjecture. According to the supplied statement-verification note, the report proves the displayed property with a gross-model qualification for NIP theories, observes that it does not characterize NIP, and conjectures only the existence of some unspecified strengthening. These are supplied descriptions of the report, not independently verified claims.
\subsection{Short English statement}
Find a stronger model-cardinality condition that holds exactly for NIP theories. The proposed strengthening, the cardinal quantifiers, the meaning of \(S(M)\), and the definition of a gross model are missing, so the supplied record does not determine a precise conjecture.
\subsection{Sources}
\begin{itemize}
\item Supplied local UnsolvedMath registry record, permanent ID \(30001406\), title \emph{Cardinality Characterization of NIP Theories}, exact source problem code \texttt{catalogue entry}, in the Oberwolfach Reports Open Problems collection. The recorded proposed year is 2010; no individual proposer is named.
\item Supplied report citation: \emph{Model Theory: Around Valued Fields and Dependent Theories} (2010), Oberwolfach Reports, pp.\ 5--53, DOI \texttt{10.4171/owr/2010/01}. Source URL and supplied literature-source link, labeled ``original Oberwolfach problem report'': \url{https://doi.org/10.4171/owr/2010/01}.
\item Supplied literature assessment checked on \texttt{2026-08-21}: the record labels the problem open and states that the report supplies no proof of the conjectured unspecified strengthening, and that no later primary resolution of that strengthening was located. This is supplied dated curation, not an independently verified literature assessment. The background explicitly requests independent verification of time-sensitive status and cited claims.
\item The supplied statement-verification note classifies the statement as \texttt{unrecoverable}, supplies no separate clean statement, and says that the original record is retained publicly for provenance. Its account of the report's gross-model result and failure of that property to characterize NIP has not been independently verified here.
\end{itemize}
\end{document}
